% GENERATED FILE. Edit canonical lessons, then run npm run generate:books.
% canonical-source-hash: fnv1a64:8bb9faeccd7a4fa8
% canonical-lessons: BN-C411-dakghor, BN-C411-ophis, BN-C411-chithi, BN-C411-piyon, BN-C411-pathano

\chapter{Post and Offices}
\label{ch:bn-post-and-offices}
\hlchaptermodality{411}

\begin{chapteropening}
\textbf{By the end of this chapter:} \emph{I can send a letter or a parcel from the post office.}

Complete the last lesson of chapter 411: I can send a letter or a parcel from the post office.
\end{chapteropening}

\section[\texorpdfstring{\bn{ডাকঘর} (ḍākghôr)}{ডাকঘর (ḍākghôr)}]{\bn{ডাকঘর} (ḍākghôr) — a post office}

\label{lesson:BN-C411-dakghor}



\begin{quote}
Before the new one: say the Bengali for a square (open place), then the Bengali for to pace up and down.
\end{quote}

\subsection*{You'll want to know: \bn{ডাকঘর}}
\textbf{\bn{ডাকঘর}} — \emph{ḍākghôr} — \textquotedblleft{}a post office\textquotedblright{}.

\begin{quote}
\bn{আমি} \bn{ডাকঘর} \bn{থেকে} \bn{একটা} \bn{খাম} \bn{কিনলাম।}
\end{quote}

\begin{quote}
\emph{āmi ḍākghôr theke ekṭā khām kinlām.}
\end{quote}

\begin{quote}
I bought an envelope at the post office.
\end{quote}

\bn{ডাক} is the post itself: \bn{ডাকে} \bn{আসা}, to come by post.

\begin{cousinweb}
\textbf{\bn{ডাকঘর}} (\emph{ḍākghôr}) joins \bn{ডাক}, which CDIAL lists as Bengali \emph{ḍāk}, \textquotedblleft{}a relay, the post, a posting stage\textquotedblright{} (CDIAL 5519), to \textbf{\bn{ঘর}}, \textquotedblleft{}a house\textquotedblright{}: the house of the post.
\end{cousinweb}

\begin{grammarlens}[title={What you've built}]
One more word for getting around.
\end{grammarlens}

\subsection*{Your turn}
\begin{itemize}
\raggedright
  \item \emph{Say it:} \emph{ḍākghôr}
  \item \emph{Say it:} \emph{ḍākghôr}, once more
  \item \emph{Say it:} \emph{āmi ḍākghôr theke ekṭā khām kinlām.} — I bought an envelope at the post office.
  \item \emph{Recall:} say the Bengali for a square (open place), then the Bengali for to pace up and down, then say \emph{ḍākghôr} again
\end{itemize}

\begin{tcolorbox}[breakable,colback=teal!4,colframe=teal!35!black,title={Before you move on}]
What does \bn{ডাকঘর} mean? (\textquotedblleft{}A post office\textquotedblright{}.) Say it once more.
\end{tcolorbox}
\section[\texorpdfstring{\bn{অফিস} (ôphis)}{অফিস (ôphis)}]{\bn{অফিস} (ôphis) — an office}

\label{lesson:BN-C411-ophis}



\begin{quote}
Before the new one: say the Bengali for to pace up and down, then the Bengali for a post office.
\end{quote}

\subsection*{You'll want to know: \bn{অফিস}}
\textbf{\bn{অফিস}} — \emph{ôphis} — \textquotedblleft{}an office\textquotedblright{}.

\begin{quote}
\bn{মীরার} \bn{অফিস} \bn{স্টেশনের} \bn{কাছে।}
\end{quote}

\begin{quote}
\emph{mīrār ôphis sṭeshôner kāchhe.}
\end{quote}

\begin{quote}
Mira's office is near the station.
\end{quote}

\begin{cousinweb}
\textbf{\bn{অফিস}} (\emph{ôphis}) is borrowed from English \emph{office}; the first letter is \bn{অ}, said as an open ô.
\end{cousinweb}

\begin{grammarlens}[title={What you've built}]
One more word for getting around.
\end{grammarlens}

\subsection*{Your turn}
\begin{itemize}
\raggedright
  \item \emph{Say it:} \emph{ôphis}
  \item \emph{Say it:} \emph{ôphis}, once more
  \item \emph{Say it:} \emph{mīrār ôphis sṭeshôner kāchhe.} — Mira's office is near the station.
  \item \emph{Recall:} say the Bengali for to pace up and down, then the Bengali for a post office, then say \emph{ôphis} again
\end{itemize}

\begin{tcolorbox}[breakable,colback=teal!4,colframe=teal!35!black,title={Before you move on}]
What does \bn{অফিস} mean? (\textquotedblleft{}An office\textquotedblright{}.) Say it once more.
\end{tcolorbox}
\section[\texorpdfstring{\bn{চিঠি} (chiṭhi)}{চিঠি (chiṭhi)}]{\bn{চিঠি} (chiṭhi) — a letter (sent by post)}

\label{lesson:BN-C411-chithi}



\begin{quote}
Before the new one: say the Bengali for a post office, then the Bengali for an office.
\end{quote}

\subsection*{You'll want to know: \bn{চিঠি}}
\textbf{\bn{চিঠি}} — \emph{chiṭhi} — \textquotedblleft{}a letter (sent by post)\textquotedblright{}.

\begin{quote}
\bn{মামা} \bn{বিদেশ} \bn{থেকে} \bn{একটা} \bn{চিঠি} \bn{লিখলেন।}
\end{quote}

\begin{quote}
\emph{māmā bidesh theke ekṭā chiṭhi likhlen.}
\end{quote}

\begin{quote}
My uncle wrote a letter from abroad.
\end{quote}

\begin{cousinweb}
\textbf{\bn{চিঠি}} (\emph{chiṭhi}) is inherited: CDIAL lists Bengali \emph{ciṭhi}, \textquotedblleft{}a letter\textquotedblright{}, under the reconstructed \emph{ciṣṭa} (CDIAL 4832).
\end{cousinweb}

\begin{grammarlens}[title={What you've built}]
One more word for getting around.
\end{grammarlens}

\subsection*{Your turn}
\begin{itemize}
\raggedright
  \item \emph{Say it:} \emph{chiṭhi}
  \item \emph{Say it:} \emph{chiṭhi}, once more
  \item \emph{Say it:} \emph{māmā bidesh theke ekṭā chiṭhi likhlen.} — My uncle wrote a letter from abroad.
  \item \emph{Recall:} say the Bengali for a post office, then the Bengali for an office, then say \emph{chiṭhi} again
\end{itemize}

\begin{tcolorbox}[breakable,colback=teal!4,colframe=teal!35!black,title={Before you move on}]
What does \bn{চিঠি} mean? (\textquotedblleft{}A letter (sent by post)\textquotedblright{}.) Say it once more.
\end{tcolorbox}
\section[\texorpdfstring{\bn{পিয়ন} (piyôn)}{পিয়ন (piyôn)}]{\bn{পিয়ন} (piyôn) — a postman; an office messenger}

\label{lesson:BN-C411-piyon}



\begin{quote}
Before the new one: say the Bengali for an office, then the Bengali for a letter (sent by post).
\end{quote}

\subsection*{You'll want to know: \bn{পিয়ন}}
\textbf{\bn{পিয়ন}} — \emph{piyôn} — \textquotedblleft{}a postman; an office messenger\textquotedblright{}.

\begin{quote}
\bn{পিয়ন} \bn{দুপুরে} \bn{আমাদের} \bn{চিঠি} \bn{দিয়ে} \bn{গেল।}
\end{quote}

\begin{quote}
\emph{piyôn dupure āmāder chiṭhi diye gelô.}
\end{quote}

\begin{quote}
The postman left our letters at midday.
\end{quote}

\begin{cousinweb}
\textbf{\bn{পিয়ন}} (\emph{piyôn}) is usually traced to Portuguese \emph{peão}, \textquotedblleft{}a foot-soldier, a man on foot\textquotedblright{}, the same word as English \emph{peon}.
\end{cousinweb}

\begin{grammarlens}[title={What you've built}]
One more word for getting around.
\end{grammarlens}

\subsection*{Your turn}
\begin{itemize}
\raggedright
  \item \emph{Say it:} \emph{piyôn}
  \item \emph{Say it:} \emph{piyôn}, once more
  \item \emph{Say it:} \emph{piyôn dupure āmāder chiṭhi diye gelô.} — The postman left our letters at midday.
  \item \emph{Recall:} say the Bengali for an office, then the Bengali for a letter (sent by post), then say \emph{piyôn} again
\end{itemize}

\begin{tcolorbox}[breakable,colback=teal!4,colframe=teal!35!black,title={Before you move on}]
What does \bn{পিয়ন} mean? (\textquotedblleft{}A postman; an office messenger\textquotedblright{}.) Say it once more.
\end{tcolorbox}
\section[\texorpdfstring{\bn{পাঠানো} (pāṭhāno)}{পাঠানো (pāṭhāno)}]{\bn{পাঠানো} (pāṭhāno) — to send}

\label{lesson:BN-C411-pathano}



\begin{quote}
Before the new one: say the Bengali for a letter (sent by post), then the Bengali for a postman.
\end{quote}

\subsection*{You'll want to know: \bn{পাঠানো}}
\textbf{\bn{পাঠানো}} — \emph{pāṭhāno} — \textquotedblleft{}to send\textquotedblright{}.

\begin{quote}
\bn{আমি} \bn{ডাকঘর} \bn{থেকে} \bn{দিদিমাকে} \bn{একটা} \bn{পার্সেল} \bn{পাঠালাম।}
\end{quote}

\begin{quote}
\emph{āmi ḍākghôr theke didimāke ekṭā pārsel pāṭhālām.}
\end{quote}

\begin{quote}
I sent my grandmother a parcel from the post office.
\end{quote}

Past (he/she …): \bn{তিনি} \bn{চিঠি} \bn{পাঠালেন}, \emph{tini chiṭhi pāṭhālen}.

The one you send to takes -\bn{কে}: \bn{বাবাকে} \bn{টাকা} \bn{পাঠানো}, to send money to Father. Also for people: \bn{ছেলেকে} \bn{স্কুলে} \bn{পাঠানো}, to send a son to school.

\begin{cousinweb}
\textbf{\bn{পাঠানো}} (\emph{pāṭhāno}) is inherited: CDIAL lists Bengali \emph{paṭhāna}, \textquotedblleft{}sends\textquotedblright{}, under Sanskrit \emph{prátiṣṭhati}, \textquotedblleft{}sets out\textquotedblright{}, and its causing form \emph{prásthāpayati}, \textquotedblleft{}sends\textquotedblright{} (CDIAL 8607): to send is to make someone set out.
\end{cousinweb}

\begin{grammarlens}[title={What you've built}]
One more word for getting around.
\end{grammarlens}

\subsection*{Your turn}
\begin{itemize}
\raggedright
  \item \emph{Say it:} \emph{pāṭhāno}
  \item \emph{Say it:} \emph{pāṭhāno}, once more
  \item \emph{Say it:} \emph{āmi ḍākghôr theke didimāke ekṭā pārsel pāṭhālām.} — I sent my grandmother a parcel from the post office.
  \item \emph{Recall:} say the Bengali for a letter (sent by post), then the Bengali for a postman, then say \emph{pāṭhāno} again
\end{itemize}

\begin{tcolorbox}[breakable,colback=teal!4,colframe=teal!35!black,title={Before you move on}]
What does \bn{পাঠানো} mean? (\textquotedblleft{}To send\textquotedblright{}.) Say it once more.
\end{tcolorbox}
